\section{Polynomial-Linked Tucker Models}

\subsection{Model and Observation Assumption}

Let \(\cY \in \R^{I_1 \times \cdots \times I_N}\) be an observed tensor and let
\begin{equation}
  \cS =
  \tucker{\cX}{\bA^{(1)}, \ldots, \bA^{(N)}},
  \qquad
  \cX \in \R^{R_1 \times \cdots \times R_N},
  \quad
  \bA^{(n)} \in \R^{I_n \times R_n}.
\end{equation}
NTD-PL treats \(\cS\) as a shared low-rank latent tensor and models each
observed entry through a scalar polynomial response:
\begin{equation}
  Y_{\mathbf i}
  =
  f_{\bbeta}(S_{\mathbf i})+\varepsilon_{\mathbf i},
  \qquad
  f_{\bbeta}(s)=\sum_{p=0}^{P}\beta_p s^p,
  \qquad
  \mathbf i=(i_1,\ldots,i_N).
\end{equation}
Equivalently, the tensor prediction is
\begin{equation}
  \widehat{\cY}
  =
  f_{\bbeta}(\cS)
  =
  \sum_{p=0}^{P} \beta_p \cS^{\odot p},
  \qquad
  \bbeta \in \R^{P+1}.
\end{equation}
The constant term \(\beta_0\) captures affine response shifts, the linear term
\(\beta_1\cS\) is the usual Tucker channel, and terms with \(p\ge2\) activate
nonlinear channels on the same latent tensor. Standard Tucker is recovered by
setting \(\beta_1=1\) and \(\beta_p=0\) for \(p\neq 1\). Thus NTD-PL changes
the entry-generation map while leaving the Tucker backbone and its
multilinear ranks fixed.

\subsection{Shared-Backbone High-Order Interactions}

Although \(f_{\bbeta}\) is scalar, the term \(\cS^{\odot p}\) is not merely an
independent output calibration. Expanding one entry gives
\begin{align}
  S_{\mathbf i}^{p}
  &=
  \left(
  \sum_{r_1=1}^{R_1}\cdots\sum_{r_N=1}^{R_N}
  \cX_{r_1,\ldots,r_N}
  \prod_{n=1}^{N} A^{(n)}_{i_n,r_n}
  \right)^p \\
  &=
  \sum_{\mathbf r^{(1)},\ldots,\mathbf r^{(p)}}
  \left(
  \prod_{q=1}^{p} \cX_{r_1^{(q)},\ldots,r_N^{(q)}}
  \right)
  \prod_{n=1}^{N}
  \prod_{q=1}^{p} A^{(n)}_{i_n,r_n^{(q)}} .
\end{align}
Thus a degree-\(p\) channel couples \(p\) latent Tucker index tuples while
reusing the same mode factors in every copy. This is the sense in which the
polynomial link acts as a shared-backbone high-order Tucker interaction:
it forms the kind of products that would appear in an explicit \(p\)-fold
Tucker expansion, but ties all mode projections to
\(\bA^{(1)},\ldots,\bA^{(N)}\) and mixes the resulting channels only through
the scalar coefficients \(\{\beta_p\}\).

This distinction also separates NTD-PL from post-hoc calibration. A post-hoc
calibrator learns a scalar map after a Tucker tensor has already been fitted.
NTD-PL instead learns the latent tensor and the link jointly, so the nonlinear
response can shape the low-rank geometry itself. Section~\ref{sec:theory}
formalizes the induced rank expansion and the resulting separation from
standard Tucker.

\subsection{Masked Least-Squares Objective}

Let \(\Omega\in\{0,1\}^{I_1\times\cdots\times I_N}\) denote the observed-entry
mask and let \(|\Omega|=\sum_{\mathbf i}\Omega_{\mathbf i}\). The empirical
data-fitting term is the average squared residual over observed entries,
\begin{equation}
  \mathcal L_{\Omega}
  =
  \frac{1}{2|\Omega|}
  \sum_{\mathbf i:\Omega_{\mathbf i}=1}
  \left(Y_{\mathbf i}-f_{\bbeta}(S_{\mathbf i})\right)^2 .
\end{equation}
The squared loss is the standard least-squares choice for real-valued
reconstruction and completion, and it aligns with the RMSE-based evaluation
used in the experiments. The Frobenius form below is simply the tensor notation
for the same observed-entry residual sum:
\begin{equation}
  \min_{\cX,\{\bA^{(n)}\},\bbeta}
  \frac{1}{2|\Omega|}
  \norm{
    \Omega \odot
    \left(\cY - f_{\bbeta}(\cS)\right)
  }_F^2
  +
  \frac{\lambda_X}{2}\norm{\cX}_F^2
  +
  \sum_{n=1}^{N}\frac{\lambda_n}{2}\norm{\bA^{(n)}}_F^2
  +
  \frac{\lambda_\beta}{2}\norm{\bbeta}_2^2.
\end{equation}
The factor \(1/|\Omega|\) keeps the loss and gradient scale comparable across
full-observation reconstruction and different missing rates, while the
factor \(1/2\) removes a constant \(2\) from the residual gradients. The
Frobenius penalties are ridge/Tikhonov regularizers for the core and mode
factors; the \(\ell_2\) penalty on \(\bbeta\) stabilizes the polynomial link
and makes the coefficient subproblem well-conditioned. When
\(\Omega\equiv\mathbf 1\), the same objective gives full tensor
reconstruction; otherwise it gives masked tensor completion.

\paragraph{Practical scaling.}
Like Tucker, NTD-PL has scale indeterminacies between the core and mode
factors. Polynomial powers make this numerically more sensitive: if the scale
of \(\cS\) drifts, different powers change at different rates and the link
coefficients become poorly conditioned. In the implementation we initialize
from Tucker, normalize factor scales during training, and absorb the inverse
scales into the core. This gauge fixing does not change the latent tensor
\(\cS\) or the model class, but it makes the alternating optimization more
stable in the low-rank regimes used in the experiments.
